\documentclass{beerydoc}


\begin{document}


\section*{\LARGE The \pkg{xgeometry} package}

\LaTeX3 implementation of the \pkg{geometry} package

Oliver Beery

Version 0.0.0\quad 8 August 2026


\section{Introduction}

\subsection{Loading the package}
\label{subsec:load}

\listheading{Requirements:}
\begin{itemize}
  \item \LaTeXe{} version 2023-11-01 or newer
  \item \pkg{l3kernel} version 2023-11-01 or newer
\end{itemize}

\subsection{Syntax}
\label{subsec:syntax}

This documentation uses the syntaxes \meta{dimen} and \meta{number}.
These syntaxes have the same meaning as the arguments to \cs{dimeval} and
\cs{fpeval}, respectively, which are documented in \pkg{usrguide}.


\section{Keys}
\label{sec:keys}

\subsection{Setting \cs{paperwidth} and \cs{paperheight}}
\label{subsec:paper}

\begin{code}{paper/width}
  \begin{syntax}
    \key{paper/width=\meta{dimen}}
  \end{syntax}
  Sets \cs{paperwidth} to \meta{dimen}.
  The initial value is \qty{21}{cm} (same as |A4|).
\end{code}

\begin{code}{paper/height}
  \begin{syntax}
    \key{paper/height=\meta{dimen}}
  \end{syntax}
  Sets \cs{paperheight} to \meta{dimen}.
  The initial value is \qty{29.7}{cm} (same as |A4|).
\end{code}

\begin{code}{paper/flipped}
  \begin{syntax}
    \key{flipped=true|false}
  \end{syntax}
  Boolean key that sets whether to switch the values of \cs{paperwidth} and
  \cs{paperheight}.
  The initial value is |false|.
\end{code}

\subsection{Setting the body}
\label{subsec:body}

\begin{code}{body/width}
  \begin{syntax}
    \key{body/width=\meta{dimen}}
  \end{syntax}
  Sets the width of the body to \meta{dimen}.
\end{code}

\begin{code}{body/height}
  \begin{syntax}
    \key{body/height=\meta{dimen}}
  \end{syntax}
  Sets the height of the body to \meta{dimen}.
\end{code}

\begin{code}{body/width-paper-ratio}
  \begin{syntax}
    \key{body/width-paper-ratio=\meta{number}}
  \end{syntax}
  Sets the ratio of the body width to \cs{paperwidth} to \meta{number}.
  The initial value is \meta{not set}.
\end{code}

\begin{code}{body/height-paper-ratio}
  \begin{syntax}
    \key{body/height-paper-ratio=\meta{number}}
  \end{syntax}
  Sets the ratio of the body height to \cs{paperheight} to \meta{number}.
  The initial value is \meta{not set}.
\end{code}

\begin{code}{body/include-header}
  \begin{syntax}
    \key{body/include-header=true|false}
  \end{syntax}
  Boolean key that set whether the header is included as part of the body.
  The initial value is |false|.
\end{code}

\begin{code}{body/include-footer}
  \begin{syntax}
    \key{body/include-footer=true|false}
  \end{syntax}
  Boolean key that set whether the footer is included as part of the body.
  The initial value is |false|.
\end{code}

\begin{code}{body/include-margin-notes}
  \begin{syntax}
    \key{body/include-margin-notes=true|false}
  \end{syntax}
  Boolean key that set whether the margin notes are included as part of the
  body.
  The initial value is |false|.
\end{code}

\subsection{Setting the margins}
\label{subsec:margins}

\begin{code}{margin/top, margin/bottom, margin/left, margin/right}
  \begin{syntax}
    \key{margin/top=\meta{dimen}}
  \end{syntax}
  Sets the top, bottom, left, or right margin to \meta{dimen}.
\end{code}

\begin{code}{margin/all}
  \begin{syntax}
    \key{margin/all=\meta{dimen}}
  \end{syntax}
  Meta key that sets the following keys to \meta{dimen}:
  \begin{itemize}
    \item \key{margin/top}
    \item \key{margin/bottom}
    \item \key{margin/left}
    \item \key{margin/right}
  \end{itemize}
\end{code}

\begin{code}{margin/top-bottom-ratio}
  \begin{syntax}
    \key{margin/top-bottom-ratio=\meta{number}}
  \end{syntax}
  Sets the ratio of the top margin to the bottom margin to \meta{number}.
  The initial value is \meta{not set}.
\end{code}

\begin{code}{margin/left-right-ratio}
  \begin{syntax}
    \key{margin/left-right-ratio=\meta{number}}
  \end{syntax}
  Sets the ratio of the left margin to the right margin to \meta{number}.
  The initial value is \meta{not set}.
\end{code}

\subsection{Resetting the keys}

\begin{code}{reset}
  Meta key that sets all the \pkg{xgeometry} package keys to their initial
  values.
\end{code}

\section{Programming}

\begin{code}{\l_xgeometry_paper_flipped_bool}
  The value of the key \key{paper/flipped} (\S\ref{subsec:paper}).
\end{code}

\begin{code}
  {
    \l_xgeometry_body_include_header_bool,
    \l_xgeometry_body_include_footer_bool,
    \l_xgeometry_body_include_margin_notes_bool
  }
  The value of the keys \key{body/include-header}, \key{body/include-footer},
  and \key{body/include-margin-notes} (\S\ref{subsec:body}).
\end{code}

\begin{code}
  {
    \xgeometry_margin_top:,
    \xgeometry_margin_bottom:,
    \xgeometry_margin_left:,
    \xgeometry_margin_right:
  }
  Expands to the length of the top, bottom, left, or right margin.
\end{code}

\begin{code}{\xgeometry_body_width:, \xgeometry_body_height:}
  Expands to the width or height of the body.
\end{code}


\end{document}
